%% ma: L12-B14-0019
%% lop: 12
%% chuong: 5
%% bai: 14
%% dang: 12.14.3
%% loai: TN4
%% muc: NB
%% dap_an: "A"
%% nguon: "CĐ2.5.pdf giáo viên gửi 27/10/2026 (M07, Lê Minh Khôi – Nguyễn Anh Khôi, Chương II Chủ đề 5 Tích có hướng của hai vectơ), Câu 3"
%% kien_thuc: "Bài 14 (tích có hướng của hai vectơ)"
%% trang_thai: chinh-thuc
\begin{ex}
Trong không gian $Oxyz$, cho hai vectơ $\vec m=(-2;6;2)$ và $\vec n=(1;-7;3)$. Tích có hướng của $\vec m$, $\vec n$ là
\choice{\True $\left[\vec m,\vec n\right]=(32;8;8)$}{$\left[\vec m,\vec n\right]=(4;1;1)$}{$\left[\vec m,\vec n\right]=(16;4;4)$}{$\left[\vec m,\vec n\right]=(-32;-8;-8)$}
\loigiai{$\left[\vec m,\vec n\right]=(18+14;2+6;14-6)=(32;8;8)$.}
\end{ex}
